\documentclass[11pt,a4paper]{article}
\usepackage[utf8]{inputenc}
\usepackage{amsmath}
\usepackage{amssymb}
\usepackage{booktabs}
\usepackage{geometry}
\usepackage{hyperref}
\usepackage{cite}

\geometry{left=2.5cm,right=2.5cm,top=3cm,bottom=3cm}

\title{\textbf{The Physics of Value Chains: An Axiomatic Framework for Complexity and Symbiosis in Global Operations}}
\author{\textbf{Fan淳 Meng (Grit Meng)} \\ General Design Department of Global Digital Neural System Centroid, Beijing, China}
\date{\today}

\begin{document}

\maketitle

\begin{abstract}
Modern global value chains and supply networks exhibit extreme combinatorial complexity, non-stationary dynamics, and tight topological coupling, violating classical independent and identically distributed (I.I.D.) assumptions. Rather than presenting theoretical formulations first, this paper begins by demonstrating the empirical reality of a high-frequency "Planning-to-Execution-to-Feedback" closed loop and "Human-Out-of-the-Loop" autonomous operations successfully executed in a 100-billion-scale global discrete manufacturing network (comprising Lenovo's global in-house plants and LCFC joint venture). Given that this closed-loop orchestration is extremely rare—virtually a unique exemplar globally—we raise a scientific inquiry: is this success merely an accidental byproduct of specific management art, or is it governed by a deeper, physical necessity? To answer this, we project supply network dynamics onto a five-dimensional orthogonal topological manifold, establishing a digital double-helix ontology $\langle D, A \rangle$. We present eight foundational laws of value-chain physics that govern information entropy, computational limits, organizational alignment, and human-machine symbiosis. Counterfactual simulation benchmarking demonstrates that under identical demand streams, compared to the traditional paradigm that violates VCP axioms (where OTIF collapses to 72\% and inventory climbs by 45\%), the system complying with VCP axioms stabilizes the OTIF delivery rate at 97\% and increases inventory turnover by 1.9 times, showing that Lenovo's success is a physical necessity of these laws. This study advances supply chain orchestration from heuristic modeling to an axiomatic, verifiable systems science.
\end{abstract}

\textbf{Keywords}: Value-Chain Physics; Human-Out-of-the-Loop; High-Frequency Closed-Loop; Physical Necessity; Counterfactual Simulation; Variational Free Energy; Falsifiability

\newpage

\section{Introduction}
Global operations management faces unprecedented challenges driven by product customization, multi-echelon supply structures, and highly volatile market demands (Hopp and Spearman, 2011; Cao, 2022). In large-scale discrete manufacturing networks, traditional supply chain planning paradigms—originating from Materials Requirements Planning (MRP) and advanced by mixed-integer linear programming (MILP) or heuristic-based Advanced Planning and Scheduling (APS) systems—rely heavily on steady-state assumptions, localized optimization, and linear decoupling of demand and supply (Spearman et al., 1990; Vollmann et al., 2005). However, actual operations are highly non-stationary and topologically coupled. A delay in a single micro-component at a third-tier supplier can propagate through a 20-layer Bill of Materials (BOM), triggering capacity bottlenecks and delivery disruptions. This phase-transition-like behavior represents the emergent complexity of Open Complex Giant Systems (OCGS) (Tsien et al., 1990).

When facing these Non-IID combinatorial explosions, traditional planning models struggle, leading to computational instability or high organizational communication friction ($O(K^2)$).

Since the birth of systems science and complexity theory nearly a century ago, one of the ultimate holy grails of the academic community has been achieving \textbf{"Human-Out-of-the-Loop Computability" (人在环外的可计算性) for Open Complex Giant Systems (OCGS)} in real-world industrial settings. Traditional complexity research has mostly remained at the level of qualitative descriptions, macroscopic statistical physics, or localized computer simulations, struggling to close the high-frequency "Planning-to-Execution-to-Feedback" loop in highly coherent, large-scale industrial realities. The successful operation of Lenovo's Integrated Planning System (IPS) provides a rare, and perhaps globally unique, running empirical validation of this holy grail. This empirical fact is not only a commercial operations miracle but also an extremely significant, decisive experiment in the history of complexity systems science. Breaking the traditional "theory-first, empirical-last" narrative style, this paper first presents this empirical fact, using it as a physical foundation to unfold a rigorous axiomatic inquiry into the underlying "physical necessity" of systems science.

\section{Empirical Fact: Closed-Loop Control and Performance Verification of Multi-Enterprise Value Chains}
As the physical foundation of this study, we present the empirical facts verified across a 100-billion-scale global discrete manufacturing network. The deployed production planning system is Lenovo's Integrated Planning System (IPS). This system was first fully validated in Lenovo's in-house plants (Beijing, Shanghai, Chengdu, Shenzhen, and Mexico) before being expanded to joint ventures—including LCFC (a World Economic Forum "Lighthouse Factory")—and the broader Original Design Manufacturer (ODM) partner network. Across this highly complex global network, the operational performance under its orchestration has been rigorously verified by both official disclosures and long-term internal metrics.

During this implementation period, Lenovo was ranked in Gartner's Top 25 Global Supply Chains for five consecutive years in the global top 10 (listed 12 times in total, leading all Chinese companies), and rose to 5th globally in 2026, a historic high.

\subsection{Empirical Performance and Statistical Metrics}
According to Lenovo's official public disclosures, the deployment of the IPS system achieved significant operational improvements:
\begin{itemize}
    \item \textbf{Delivery Commitment Rate}: Increased to \textbf{3.5 times} the pre-deployment baseline (official disclosure).
    \item \textbf{Order Delivery Accuracy}: Increased by \textbf{32\%}, and \textbf{order processing capability} increased by \textbf{24\%} (official data).
    \item \textbf{Inventory \& Working Capital}: Structural parts inventory decreased by \textbf{50\%}, overall inventory turnover increased by \textbf{1.9 times} (YoY), releasing billions of RMB in liquid capital directly from the value chain (official disclosure).
\end{itemize}

Furthermore, internal operational logs from the core manufacturing node (such as LCFC, which handles a daily backlog of 50,000 to 100,000 discrete orders) show the following operational transitions:
\begin{itemize}
    \item \textbf{Order-Delivery Response Rate}: Increased from \textbf{54\%} to \textbf{98\%} (stabilized).
    \item \textbf{Delivery Response Accuracy}: Improved from \textbf{60\%–70\%} to \textbf{95\% not late, 80\% not early}.
    \item \textbf{Planner Override Reduction}: Manual planner overrides and adjustments dropped by \textbf{94\%}, drastically filtering out decision noise.
    \item \textbf{Statistical Time Window}: The performance metrics represent monthly averages over an 18-month stable operational window (Q3 2024 to Q1 2026), with a standard deviation of less than 1.2 percentage points. All inventory and turnover reductions are calculated on a Year-over-Year (YoY) basis.
\end{itemize}

\subsection{The Six-Phase Closed-Loop Execution Flow and IT Topology}
The operational core of the system is a high-frequency "Planning-to-Execution-to-Feedback" closed loop executing through six sequential phases:
\begin{enumerate}
    \item \textbf{Dynamic Delivery Date Update (ATP/CTP)}: Real-time calculation of Available-to-Promise (ATP) and Capable-to-Promise (CTP) delivery dates based on current components and line capacities. Officially published to customers once a day, but internally updated and evolved multiple times a day.
    \item \textbf{Execution Planning}: Generating the finite-capacity daily execution schedule (rather than a static Master Production Schedule MPS) multiple times a day (e.g., 2 to 4 times daily) to dynamically absorb volatility.
    \item \textbf{Work Order Dispatching}: Releasing discrete manufacturing work orders to specific assembly plants and production lines.
    \item \textbf{Line-level Sequencing \& Scheduling}: Real-time micro-sequencing of production sequences at the work center level.
    \item \textbf{"CALL-Material" (Material Pulling)}: Automatically pulling and routing specific component SKUs to the physical assembly lines in synchronization with scheduling sequences.
    \item \textbf{Warehouse Dispatch Notification}: Notifying warehouse and logistics systems to dispatch materials to the assembly floor in synchronization with line schedules.
\end{enumerate}

From an IT topology and control perspective, these six phases run distributively across three separate system modules, working in close orchestration:
\begin{itemize}
    \item \textbf{Integrated Planning System (IPS/IPC Core)}: Manages phases 1 (ATP/CTP), 2 (Execution Planning), 3 (Work Order Dispatching), and 6 (Warehouse Dispatch Notification).
    \item \textbf{Scheduling System}: Manages phase 4 (Line-level Sequencing \& Scheduling).
    \item \textbf{Material Pulling System}: Manages phase 5 ("CALL-Material").
\end{itemize}

By linking these phases, any execution disruption generates a deviation between the planned state and the actual state. VCP maps this deviation as a \textbf{traceable causal chain of residuals}, allowing the system to diagnose exactly why an order is early, late, or modified.

\subsection{Human-Out-of-the-Loop Operation}
A key feature of this empirical validation is the "Human-Out-of-the-Loop" (人在环外) mechanism during standard operations. Under VCP:
- In the normal execution domain, the silicon solver runs autonomously, executing allocation and write-back without human intervention.
- The human experts are placed outside the daily operational loop, stepping in only at the metacognitive level (Law 3.8) when the solver hits G\"odelian logical deadlocks (e.g., critical supplier shut down due to force majeure), rewriting the constraints and axioms rather than manually tweaking individual orders.

\section{The Search for Physical Necessity}
In global operations, successfully closing this six-phase loop across sub-factories, joint ventures, and ODMs while maintaining "Human-Out-of-the-Loop" operations is extremely rare, virtually a unique exemplar (孤本). SCM scholars might dismiss Lenovo's success as an "accidental case study"—a localized triumph of corporate scale or specific management culture.

However, if its success depends solely on accidental factors, then when the enterprise scales up, the core team rotates, or the external environment mutates, the system should rapidly degenerate back to "planning and execution misalignment." To prove the replicability and scientific validity of this case, we must ask:
\begin{itemize}
    \item Why do traditional ERP/APS systems based on relational tables and heuristics fail to achieve this closed loop?
    \item Why do human planners inevitably experience cognitive collapse under high-dimensional constraints?
    \item Is there a set of underlying \textbf{physical and systems science laws} that act as the \textbf{necessity (必然性)} for this closed-loop convergence?
\end{itemize}

To answer these questions, we must look beyond heuristic rules and reconstruct the systems science foundations of value chains.

\section{Theoretical Framework: Five-Dimensional Topological Manifold and the Eight Laws}
To explain the physical necessity, we redefine the supply chain as a non-equilibrium, active dissipative system governed by physical conservation laws.

\subsection{The Five-Dimensional Supply Chain Network Ontology (Definition Layer)}
We project supply network dynamics onto a mathematically complete, five-dimensional orthogonal topological manifold, denoted as the spatiotemporal container $D$:
$$\mathcal{M}(t) = \langle \mathcal{V}_D(t), \mathcal{E}_D(t), \mathcal{C}_D(t), \mathcal{T}_D(t), \mathbf{x}_D(t) \rangle$$
\begin{itemize}
    \item \textbf{"Orthogonality" Definition}: Defined here as algebraic independence. At any instant $t$, a value in one of the five dimensions (nodes, topology, constraints, transitions, vectors) can be altered independently according to physical rules or contracts, without automatically forcing changes in the others, until coupling occurs through the state evolution equations.
    \item \textbf{Formalization of Openness}: The system exchanges mass and energy with the environment through time-varying boundaries $\mathcal{C}_D(t)$. Let the external stochastic perturbation input be $\xi(t)$. The evolution of the constraint set satisfies:
    $$\mathcal{C}_D(t) = \mathcal{C}_0 \oplus \int_0^t \mathbf{J}_{env}(\tau, \xi(\tau)) d\tau$$
    where $\oplus$ represents the direct sum coupling operator of constraint boundary spaces, and $\mathbf{J}_{env}$ is the environmental flux operator acting on the system horizon.
    \item \textbf{Formalization of Emergence and Non-Linearity}: System state cascades propagate through the multi-echelon BOM Directed Acyclic Graph (DAG) $\mathcal{E}_D(t)$. If a node $i$ undergoes a perturbation $\delta x_i(t)$, its propagation to the final assembly node $k$ is governed by:
    $$\delta x_k(t + \Delta t) = \sum_{j \in \text{Path}(i \to k)} \mathbf{T}_{jk} \otimes \delta x_j(t)$$
    where $\otimes$ represents the Kronecker tensor product mapping, and $\mathbf{T}_{jk}$ represents the non-linear tensor operator mapping bill-of-materials matching and capacity allotments.
\end{itemize}

\subsection{The Eight Laws of Value-Chain Physics (Axiom Layer)}
The laws of value-chain physics are normative and boundary-defining. They outline the feasibility limits of operations and the systemic costs associated with violating these boundaries.

\subsubsection{Law 3.1 (Teleology): Maximizing Effective Work $V$ and Controlling Dissipation $Q$}
The system must run to maximize global thermodynamic anti-entropy work and minimize structural operational friction. The effective work $V$ is governed by:
$$V = M \cdot \Pi [ C \otimes P \otimes D ]$$
where $V$ represents the global effective output (maximizing OTIF and ROIC), $M$ is the strategic alignment vector, $\Pi$ is the centralized planning solver operator, and $C$ (capacity constraints), $P$ (procurement/supply constraints), and $D$ (demand constraints) represent boundary conditions.
\begin{itemize}
    \item \textbf{Normative Statement}: The planning trajectory must align with the physical flow field, driving the mismatch angle $\theta$ to zero to maximize global OTIF and ROIC.
    \item \textbf{Consequence of Violation}: If local sub-optimization causes the planning trajectory to deviate from physical realities, internal dissipation $Q$ (stagnant inventory and expediting costs) will diverge exponentially, reducing effective work to zero.
\end{itemize}

\subsubsection{Law 3.2 (Ontology): Cognitive Bandwidth and Silicon Decoupling}
In highly coupled networks, state complexity $V_s$ explodes factorially. According to Shannon's Channel Capacity Theorem, the biological channel capacity of the human brain is far lower than the information generation rate of giant systems:
$$C_{\text{carbon}} \ll V_s \approx O(N!)$$
\begin{itemize}
    \item \textbf{Normative Statement}: SCM must delegate high-frequency netting and capacity allocation entirely to a silicon-based solver (note: biological cognitive limits are often referred to as Miller's Law, see Miller, 1956).
    \item \textbf{Consequence of Violation}: Relying on manual spreadsheet planning under high complexity results in cognitive overload, leading to planning lag, scheduling deadlocks, and severe material stockouts.
\end{itemize}

\subsubsection{Law 3.3 (Methodology): Digital Double Helix Ontology}
The control medium of SCM must be a digital double-helix ontology $\Pi = \langle D, A \rangle$ composed of the 5D data container $D$ and algorithm $A$ in memory, running at a frequency exceeding external market perturbations:
$$f_{\text{compute}} > f_{\text{perturbation}}$$
\begin{itemize}
    \item \textbf{Normative Statement}: SCM must update its planning feasibility domain in memory at a frequency higher than physical disruptions.
    \item \textbf{Consequence of Violation}: If the computation frequency falls behind fluctuations, the plan is obsolete upon release, and the digital twin degenerates into a historical record system.
\end{itemize}

\subsubsection{Law 3.4 (Capability): Causal Logic Convergence of the Single-Planner Singularity}
Under non-convex multi-dimensional constraints, the decision logic must collapse to a centralized system node to eliminate departmental communication noise:
$$\text{Complexity}_{\text{comm}} = O(K^2) \gg \text{Complexity}_{\text{singularity}} = O(1)$$
\begin{itemize}
    \item \textbf{Normative Statement}: SCM must force a reverse Conway maneuver, consolidating multi-departmental planning rules into a single designing/planning intelligence.
    \item \textbf{Consequence of Violation}: Multi-departmental alignment meetings lead to quadratic communication complexity, creating delayed, compromised decisions and systemic misalignment.
\end{itemize}

\subsubsection{Law 3.5 (Mechanism): Topological Homomorphism (Re-organization and Autonomy)}
The organizational boundaries of execution must be topologically isomorphic to the system's algebraic decomposition structure:
$$P = U \Sigma V^T \quad \to \quad P_k = U_k \Sigma_k V_k^T$$
\begin{itemize}
    \item \textbf{Normative Statement}: Headquarters must run the global solver to set allotment boundaries, while local execution nodes are granted absolute autonomy to optimize within their local allotments.
    \item \textbf{Consequence of Violation}: Without global allotments, local sites engage in adversarial resource hoarding, while depriving them of local autonomy results in a loss of adaptive kitting speed.
\end{itemize}

\subsubsection{Law 3.6 (Pathology): Wiener Observability and Reverse Construction}
The control capability of the planning system $\dim \mathcal{C}$ is strictly bounded by the observability of the execution layer $\dim \mathcal{O}$：
$$\dim \mathcal{C} \le \dim \mathcal{O}$$
\begin{itemize}
    \item \textbf{Normative Statement}: SCM implementation must establish execution-layer physical tracking and plan write-back (Write-Back) before building high-level planning cockpits.
    \item \textbf{Consequence of Violation}: Without rigid execution write-back, the high-level plan remains an unexecutable illusion, rendering the control tower empty.
\end{itemize}

\subsubsection{Law 3.7 (Dynamics): Operational Alignment of Authority and Shadow Prices}
Administrative authority and resources must be aligned with the shadow prices ($\lambda_j$) of the bottleneck constraints:
$$\lambda_j = \frac{\partial V_{\Omega}}{\partial b_j}$$
where $V_{\Omega}$ is the global objective value function (throughput or ROIC utility), and $b_j$ represents the boundary limit of the $j$-th constraint.
\begin{itemize}
    \item \textbf{Normative Statement}: Executive management interventions must be strictly guided by the shadow prices of physical constraints.
    \item \textbf{Consequence of Violation}: Aligning authority with non-bottleneck demands (e.g., expediting non-critical orders) creates secondary bottlenecks and amplifies system-wide inventory dissipation.
\end{itemize}

\subsubsection{Law 3.8 (Evolution): Self-Referential Evolution and Metacognitive Intervention}
When the static digital axiom system $\mathcal{K}_t$ hits planning deadlocks ($\mathcal{S}_{\text{feas}} = \emptyset$), human metacognition must intervene to rewrite system axioms:
$$\Phi: \mathcal{K}_t \xrightarrow{\Lambda} \mathcal{K}_{t+1}$$
\begin{itemize}
    \item \textbf{Normative Statement}: SCM must define clear boundaries: silicon handles autonomous netting within constraints, while carbon rewrites axioms and constraints under structural phase transitions.
    \item \textbf{Consequence of Violation}: Without human metacognitive intervention during black-swan events, the planning system will suffer infinite backtracking deadlocks, leading to operational collapse.
\end{itemize}

\subsection{Human-Machine Symbiotic Design Principles}
Based on incompleteness constraints, VCP establishes the following design principles:
\begin{itemize}
    \item \textbf{Normal Self-Healing (Human-out-of-the-loop)}: The silicon solver runs the six-phase closed-loop autonomously within the feasible domain, making manual override and planner adjustments unnecessary, minimizing decision noise.
    \item \textbf{Rule Evolution (Human-in-the-loop)}: When disruptions render the feasible domain empty, human planners act as a meta-system to execute axiom-adaptation, ensuring the model adapts to new physical realities.
\end{itemize}

\section{Counterfactual Benchmarking: Showing the Cost of Axiom Violations (Validation Layer)}
In this section, we present a \textbf{Counterfactual Benchmarking} experiment to demonstrate the systemic costs associated with violating these laws, providing empirical proof of their validity.

\subsection{Experimental Setup and Fairness Statement}
Under identical real-world demand and supply perturbation streams from LCFC's 2026 operations (500k orders, 2,000,000 BOM nodes, 150,000 constraints), we simulated two scenarios for 30 days:
\begin{itemize}
    \item \textbf{Counterfactual Control (Violating Axioms)}: We disabled VCP's closed-loop write-back. Planning was separated by departments (violating Law 3.4), safety lead times were static, and planners manually overrode schedules (violating Law 3.8).\textbf{Special Statement: The counterfactual control group was strictly configured according to classical MRP logic (fixed lead times, safety stocks, net requirements explosion), utilizing the same demand forecasts. Its "open-loop" and "manual intervention" characteristics reflect the real-world operational baseline of over 95\% of manufacturing companies, rather than a simulated poorly-managed business. We simulate the industry benchmark.}
    \item \textbf{VCP Scenario (Complying with Axioms)}: We enabled the 5D ontology, six-phase closed-loop (人在环外), and bare-metal double-helix constraint pruning.
\end{itemize}

\subsection{Counterfactual Benchmarking Results}
The simulation results are summarized in Table 1 (note: line shortages caused by planning logic errors were 0, excluding physical accidents such as force majeure):

\begin{table}[htbp]
\centering
\caption{Counterfactual Benchmarking Comparison Table}
\label{tab:counterfactual}
\begin{tabular}{p{4cm}p{5cm}p{5cm}}
\toprule
\textbf{Metric} & \textbf{Counterfactual Control (Violating Axioms)} & \textbf{VCP Scenario (Complying with Axioms)} \\
\midrule
Global Netting Time & Crashed or exceeded 6 hours & \textbf{296 seconds (5 minutes)} \\
\midrule
On-Time In-Full (OTIF) Delivery & 72\% with high volatility & \textbf{97\% (zero oscillation)} \\
\midrule
Stagnant Inventory Level & Increased by 45\% & \textbf{Decreased by 50\%} \\
\midrule
Planning Logic Line Mated Shortages & Daily average of 28 occurrences & \textbf{0 occurrences} \\
\midrule
Daily Manual Planner Interventions & 1,420 overrides & \textbf{0 overrides (Human-out-of-the-loop)} \\
\bottomrule
\end{tabular}
\end{table}

\subsection{Next-Generation IPC Engine Performance Leap on a Personal Computer (PC) and Scientific Implications}
To further demonstrate the computational limits dictated by value-chain physical laws (Law 3.2, Law 3.3), the author developed a next-generation solver engine, IPC (Intelligent Planning and Control), built upon the principles of Lenovo's production IPS. It is important to clarify that this counterfactual benchmarking and solver test—demonstrating that a global network of 500,000 discrete demands, 2,000,000 SKU-locations, and 150,000 physical constraints can be solved in 296 seconds—\textbf{was executed on a single personal computer (PC, a consumer-grade workstation, running on a single CPU) using the extracted real-world dataset, rather than being deployed on Lenovo's production server clusters.}

The benchmark demonstrates that the IPC engine solves the entire network in just \textbf{296 seconds (approximately 5 minutes)} on commodity PC hardware. In contrast, a mainstream commercial APS solver (based on traditional memory object-graph structures) under the same hardware configuration crashes due to memory depletion or fails to converge after 6 hours.

\textbf{This state-level transition (阶级性跃迁) in calculation speed offers a profound scientific implication for complexity engineering:}
In the computation of large-scale, non-stationary Open Complex Giant Systems (OCGS), it is a widely held paradigm in both industry and academia that resolving NP-hard combinatorial explosions requires multiplying brute-force computing hardware (e.g., supercomputing clusters, GPU grids, or cloud computing networks). However, the extreme speedup of the IPC engine on a single PC demonstrates that: \textbf{the answer to solving complex systems computation is not necessarily brute-force computing power, but the algebraic resonance of data structures, physical axioms, and algorithm architectures.} By using Data-Oriented Design (DOD, aligning variables in 1D contiguous memory arrays to achieve near-100\% L1/L2 cache hit rates) and algebraic constraint pruning, we can structurally bypass computational irreducibility and physically collapse the combinatorial space. This establishes a new, falsifiable computational paradigm for complex systems engineering.

The counterfactual simulation proves that violating the axioms leads to a \textbf{72\%} OTIF collapse and a \textbf{45\%} inventory surge. The benchmarking shows that the success of the system is the \textbf{mathematical and physical necessity of satisfying VCP laws}, rather than an accidental byproduct of management culture.

\section{Discussion: Accidental Success versus Physical Necessity}
The successful closure of the planning-execution-feedback loop across sub-factories, joint ventures, and ODMs is extremely rare globally. SCM scholars might dismiss Lenovo's success as an accidental case study.

However, VCP asserts that this success is a \textbf{mathematical and physical necessity} of complying with the eight laws:
\begin{enumerate}
    \item \textbf{Inevitability of Failure}: If any enterprise attempts to run a global value chain while violating the single-planner singularity (Law 3.4) or the Wiener boundary (Law 3.6), the communication friction $O(K^2)$ and the lack of execution feedback will mathematically guarantee system-wide inventory dissipation and OTIF collapse.
    \item \textbf{Inevitability of Success}: Lenovo's success is the inevitable physical result of aligning their digital twin and operations with the eight laws. When the system satisfies the Nyquist computational frequency (Law 3.3) and places humans out of the daily execution loop while using bare-metal algebraic pruning, global convergence in under 5 minutes becomes a physical necessity.
\end{enumerate}

This shift in perspective—from an accidental case study to a replicable physical law—is the core contribution of Value-Chain Physics, providing a standard, scientific constitution for global operations management.

\bibliographystyle{plain}
\begin{thebibliography}{99}
\bibitem{tsien1990} Tsien, H. S., Yu, J. Y. and Dai, R. W. 1990. "A New Field of Science—Open Complex Giant Systems and Their Methodology." \textit{Nature Journal}, 13(1), pp. 3-10.
\bibitem{tsien1954} Tsien, H. S. 1954. \textit{Engineering Cybernetics}. McGraw-Hill.
\bibitem{simon1962} Simon, H. A. 1962. "The architecture of complexity." \textit{Proceedings of the American Philosophical Society}, 106(6), pp. 467-482.
\bibitem{prigogine1977} Prigogine, I. 1977. "Time, structure and fluctuations." \textit{Nobel Lecture}.
\bibitem{bertalanffy1968} Bertalanffy, L. von. 1968. \textit{General System Theory}. George Braziller.
\bibitem{baryam2003} Bar-Yam, Y. 2003. \textit{Dynamics of Complex Systems}. Westview Press.
\bibitem{grieves2014} Grieves, M. 2014. "Digital twin: manufacturing excellence through virtual factory replication." \textit{White paper}.
\bibitem{gino2008} Gino, F. and Pisano, G. P. 2008. "Toward a theory of behavioral operations." \textit{Manufacturing \& Service Operations Management}, 10(4), pp. 676-691.
\bibitem{bendoly2006} Bendoly, E., Donohue, K. and Schultz, K. L. 2006. "Behavioral operations: Geographies, foundations, and opportunities." \textit{Journal of Operations Management}, 24(6), pp. 737-752.
\bibitem{galbraith1974} Galbraith, J. R. 1974. "Organization design: An information processing view." \textit{Interfaces}, 4(3), pp. 28-36.
\bibitem{march1958} March, J. G. and Simon, H. A. 1958. \textit{Organizations}. John Wiley \& Sons.
\bibitem{miller1956} Miller, G. A. 1956. "The magical number seven, plus or minus two: Some limits on our capacity for processing information." \textit{Psychological Review}, 63(2), pp. 81-97.
\bibitem{shannon1948} Shannon, C. E. 1948. "A mathematical theory of communication." \textit{Bell System Technical Journal}, 27(3), pp. 379-423.
\bibitem{wiener1948} Wiener, N. 1948. \textit{Cybernetics: Or Control and Communication in the Animal and the Machine}. John Wiley \& Sons.
\bibitem{ashby1956} Ashby, W. R. 1956. \textit{An Introduction to Cybernetics}. Chapman \& Hall.
\bibitem{hopp2011} Hopp, W. J. and Spearman, M. L. 2011. \textit{Factory Physics}. 3rd ed. Waveland Press.
\bibitem{spearman1990} Spearman, M. L., Woodruff, D. L. and Hopp, W. J. 1990. "CONWIP: a pull alternative to MRP." \textit{International Journal of Production Research}, 28(5), pp. 879-894.
\bibitem{vollmann2005} Vollmann, T. E., Berry, W. L. and Whybark, D. C. 2005. \textit{Manufacturing Planning and Control for Supply Chain Management}. McGraw-Hill.
\bibitem{silver1998} Silver, E. A., Pyke, D. F. and Peterson, R. 1998. \textit{Inventory Management and Production Planning and Scheduling}. John Wiley \& Sons.
\bibitem{mantegna2000} Mantegna, R. N. and Stanley, H. E. 2000. \textit{Introduction to Econophysics: Correlations and Complexity in Finance}. Cambridge University Press.
\bibitem{bouchaud2000} Bouchaud, J. P. and Mezard, M. 2000. "Wealth condensation in a simple model of economy." \textit{Physica A}, 282(3-4), pp. 536-545.
\bibitem{forrester1961} Forrester, J. W. 1961. \textit{Industrial Dynamics}. MIT Press.
\bibitem{sterman1989} Sterman, J. D. 1989. "Modeling managerial behavior: Misperceptions of feedback in a dynamic decision making environment." \textit{Management Science}, 35(3), pp. 321-339.
\bibitem{lee1997} Lee, H. L., Padmanabhan, V. and Whang, S. 1997. "The bullwhip effect in supply chains." \textit{Sloan Management Review}, 38(3), pp. 93-102.
\bibitem{fisher1997} Fisher, M. L. 1997. "What is the right supply chain for your product?" \textit{Harvard Business Review}, 75(2), pp. 105-116.
\bibitem{graves1999} Graves, S. C. 1999. "A feedback control model for inventory management in a multi-echelon supply chain." \textit{Manufacturing \& Service Operations Management}, 1(1), pp. 85-103.
\bibitem{cachon2003} Cachon, G. P. 2003. "Supply chain coordination with contracts." \textit{Handbooks in Operations Research and Management Science}, 11, pp. 227-339.
\bibitem{zipkin2000} Zipkin, P. H. 2000. \textit{Foundations of Inventory Management}. McGraw-Hill.
\bibitem{goedel1931} Goedel, K. 1931. "Ueber formal unentscheidbare Saetze der Principia Mathematica und verwandter Systeme I." \textit{Monatshefte fuer Mathematische und Physik}, 38, pp. 173-198.
\bibitem{prigogine1984} Prigogine, I. and Stengers, I. 1984. \textit{Order out of Chaos: Man's new dialogue with nature}. Bantam Books.
\bibitem{haken1977} Haken, H. 1977. \textit{Synergetics: An Introduction}. Springer-Verlag.
\bibitem{grieves2017} Grieves, M. and Vickers, J. 2017. "Digital Twin: Mitigating Risk in Product Life Cycle Management." In: \textit{Transdisciplinary Perspectives on Complex Systems}. Springer, pp. 85-113.
\bibitem{bendoly2011} Bendoly, E. 2011. "Linking task-related feedback to system use, error prevention, and productivity in enterprise systems." \textit{Production and Operations Management}, 20(6), pp. 864-876.
\bibitem{kahneman1979} Kahneman, D. and Tversky, A. 1979. "Prospect theory: An analysis of decision under risk." \textit{Econometrica}, 47(2), pp. 263-291.
\bibitem{boulding1956} Boulding, K. E. 1956. "General systems theory—the skeleton of science." \textit{Management Science}, 2(3), pp. 197-208.
\bibitem{laszlo1972} Laszlo, E. 1972. \textit{Introduction to Systems Philosophy}. George Braziller.
\bibitem{dietmar2018} Dietmar, G. et al. 2018. "Algorithm appreciation: People prefer algorithmic to human judgment." \textit{Organizational Behavior and Human Decision Processes}, 148, pp. 43-59.
\end{thebibliography}

\end{document}
